\documentclass[12pt,a4paper]{article}
\usepackage{amsmath,amsthm,amssymb,amsfonts}
\usepackage[T1]{fontenc}
\usepackage[utf8]{inputenc}
\usepackage{hyperref}
\usepackage{geometry}
\geometry{margin=2.5cm}

\newtheorem{theorem}{Theorem}[section]
\newtheorem{lemma}[theorem]{Lemma}
\newtheorem{proposition}[theorem]{Proposition}
\newtheorem{corollary}[theorem]{Corollary}
\newtheorem{definition}[theorem]{Definition}
\newtheorem{remark}[theorem]{Remark}
\newtheorem{conjecture}[theorem]{Conjecture}

\begin{document}

\title{The Collatz Conjecture and Ergodic Theory:\\
Dynamical Systems, Invariant Measures, and Bedford--McMullen Sets}
\author{Michael Fuhrmann}
\date{March 2026}
\maketitle

\begin{abstract}
This paper explores the Collatz conjecture through the lens of ergodic theory and
topological dynamics. We develop the Collatz dynamical system on the 2-adic integers,
study invariant measures and ergodicity, present Furstenberg's correspondence principle
as a bridge between combinatorial and measure-theoretic approaches, and analyse the
connection to Bedford--McMullen measures and stochastic matrix theory. The ergodic
viewpoint reveals why almost all (in Haar measure) orbits are well-behaved, while
the positive integers -- a measure-zero subset -- may or may not all converge to 1.
We also develop the connection between self-similar measures on $\{0,1\}^\mathbb{N}$
and the structure of Collatz orbits.
\end{abstract}

\tableofcontents

\newpage

\section{Introduction}

Ergodic theory studies the long-time average behaviour of dynamical systems. For the
Collatz problem, the relevant dynamical system is the Syracuse function
$S\colon \mathbb{Z}_2^{\text{odd}} \to \mathbb{Z}_2^{\text{odd}}$ acting on the 2-adic
integers. From the ergodic viewpoint, the question becomes: what is the long-run
behaviour of a typical orbit?

The surprising answer, as we shall see, is that \emph{ergodically} the situation is
well understood: almost every orbit (with respect to Haar measure) converges to 0 in
the 2-adic metric, corresponding to the orbit eventually becoming divisible by arbitrarily
high powers of 2. The difficulty of the Collatz conjecture arises precisely because
the positive integers form a measure-zero subset of $\mathbb{Z}_2$, and ergodic theorems
say nothing about measure-zero sets.

\subsection{Setting and notation}

We work with:
\begin{itemize}
  \item $\mathbb{Z}_2$: the 2-adic integers, with Haar measure $\mu$.
  \item $S\colon \mathbb{Z}_2^{\text{odd}} \to \mathbb{Z}_2^{\text{odd}}$: the Syracuse function.
  \item $T\colon \mathbb{Z}_2 \to \mathbb{Z}_2$: the Collatz function.
  \item $\sigma\colon \{0,1\}^\mathbb{N} \to \{0,1\}^\mathbb{N}$: the left shift on binary sequences.
\end{itemize}

\section{The Collatz Dynamical System}

\begin{definition}[Dynamical system]\label{def:dynsys}
A \emph{dynamical system} (or measure-preserving system) is a quadruple
$(X, \mathcal{B}, \mu, T)$ where $(X, \mathcal{B}, \mu)$ is a probability space and
$T\colon X \to X$ is a measurable map with $\mu(T^{-1}A) = \mu(A)$ for all $A \in \mathcal{B}$.
\end{definition}

\begin{definition}[Collatz dynamical system]\label{def:collatz_ds}
The \emph{Collatz dynamical system} is $(X, \mathcal{B}, \mu, T)$ where:
\begin{itemize}
  \item $X = \mathbb{Z}_2$ (the 2-adic integers),
  \item $\mathcal{B}$ is the Borel $\sigma$-algebra of $\mathbb{Z}_2$,
  \item $\mu$ is the Haar probability measure,
  \item $T\colon \mathbb{Z}_2 \to \mathbb{Z}_2$ is the Collatz function.
\end{itemize}
As established in Theorem~\ref{thm:not_measure_preserving}, this is \emph{not}
measure-preserving in general.
\end{definition}

\begin{theorem}[$T$ is not measure-preserving]\label{thm:not_measure_preserving}
The Collatz function $T\colon \mathbb{Z}_2 \to \mathbb{Z}_2$ does not preserve the
Haar measure $\mu$. Specifically, for the even integers $E = 2\mathbb{Z}_2$:
\[
  \mu(T^{-1}(E)) \;\neq\; \mu(E) = 1/2.
\]
\end{theorem}

\begin{proof}
We compute $\mu(T^{-1}(E))$ where $E = 2\mathbb{Z}_2$ is the set of even elements,
i.e.\ we need to find all $n \in \mathbb{Z}_2$ such that $T(n)$ is even.\\
\emph{Even branch:} $T(n) = n/2$ is even iff $n/2 \in 2\mathbb{Z}_2$, i.e.\ $n \in 4\mathbb{Z}_2$.
This contributes $\mu(4\mathbb{Z}_2) = 1/4$.\\
\emph{Odd branch:} $T(n) = 3n+1$ for odd $n$. Since $n$ is odd, $3n$ is odd, so $3n+1$ is even.
Thus every odd $n$ maps to an even integer, contributing $\mu(\mathbb{Z}_2^{\text{odd}}) = 1/2$.\\
Therefore $\mu(T^{-1}(E)) = 1/4 + 1/2 = 3/4 \neq 1/2 = \mu(E)$, so $T$ does not preserve $\mu$.
\end{proof}

\begin{definition}[Invariant measure]\label{def:invariant}
A measure $\nu$ on $(\mathbb{Z}_2, \mathcal{B})$ is \emph{$T$-invariant} if
$\nu(T^{-1}A) = \nu(A)$ for all $A \in \mathcal{B}$. The existence of a $T$-invariant
probability measure is guaranteed by the Krylov--Bogolyubov theorem when $T$ is continuous
(as it is here).
\end{definition}

\begin{proposition}[Existence of invariant measure]\label{prop:exist_inv}
There exists at least one $T$-invariant Borel probability measure on $\mathbb{Z}_2$.
Such a measure is supported on the $\omega$-limit set of $T$, which contains (at least)
the cycle $\{0, 0, \ldots\}$ and $\{-1/2, -1/2, \ldots\}$ (the fixed points).
\end{proposition}

\begin{proof}
By the Krylov--Bogolyubov theorem: for any $x \in \mathbb{Z}_2$, the Cesàro means
$\nu_N = (1/N)\sum_{k=0}^{N-1} \delta_{T^k(x)}$ form a tight sequence of probability
measures (by compactness of $\mathbb{Z}_2$). Any weak limit is $T$-invariant.
\end{proof}

\section{Ergodicity and Mixing Properties}

\begin{definition}[Ergodicity]\label{def:ergodic}
A measure-preserving system $(X, \mathcal{B}, \mu, T)$ is \emph{ergodic} if every
$T$-invariant measurable set $A$ (i.e.\ $T^{-1}A = A$ up to null sets) satisfies
$\mu(A) \in \{0, 1\}$.
\end{definition}

\begin{theorem}[Ergodic theorem for the Syracuse map]\label{thm:ergodic_theorem}
Let $\nu$ be the unique $S$-invariant measure on $\mathbb{Z}_2^{\text{odd}}$ that is
absolutely continuous with respect to Haar measure (whose existence follows from the
unique ergodicity of the shift). Then for $\nu$-almost every odd $x \in \mathbb{Z}_2$
and every continuous function $\phi\colon \mathbb{Z}_2^{\text{odd}} \to \mathbb{R}$:
\[
  \lim_{N \to \infty} \frac{1}{N} \sum_{k=0}^{N-1} \phi(S^k(x)) \;=\; \int_{\mathbb{Z}_2^{\text{odd}}} \phi \, d\nu.
\]
\end{theorem}

\begin{proof}
This is Birkhoff's ergodic theorem applied to the ergodic system $(\mathbb{Z}_2^{\text{odd}}, \nu, S)$.
The key step is verifying ergodicity, which follows from the mixing property of $S$
(Theorem~\ref{thm:mixing} below).
\end{proof}

\begin{theorem}[Mixing property of $S$]\label{thm:mixing}
The Syracuse map $S\colon (\mathbb{Z}_2^{\text{odd}}, \nu) \to (\mathbb{Z}_2^{\text{odd}}, \nu)$
is \emph{mixing}: for all measurable $A, B$,
\[
  \lim_{k \to \infty} \nu(S^{-k}A \cap B) \;=\; \nu(A) \cdot \nu(B).
\]
\end{theorem}

\begin{proof}[Proof sketch]
Mixing is equivalent to the spectral property that the only eigenvalue of the unitary
operator $U_S\colon L^2(\nu) \to L^2(\nu)$, $U_S f = f \circ S$, with modulus 1 is
the constant function. This follows from the exponential decay of correlations for $S$,
which is a consequence of the Fourier estimate in Lemma~\ref{lem:decay_corr}.
\end{proof}

\begin{lemma}[Exponential decay of correlations]\label{lem:decay_corr}
For all $A, B \in \mathcal{B}(\mathbb{Z}_2^{\text{odd}})$ that are determined by the
first $r$ bits of the 2-adic expansion (i.e.\ cylinder sets of depth $r$), there exist
$C > 0$ and $\rho < 1$ such that
\[
  \bigl|\nu(S^{-k}A \cap B) - \nu(A)\nu(B)\bigr| \;\leq\; C \rho^k
\]
for all $k \geq 0$.
\end{lemma}

\begin{proof}
The correlations are controlled by the spectral radius of the transfer operator
$\mathcal{L}_S\colon L^2(\nu) \to L^2(\nu)$, defined by
$\mathcal{L}_S f(x) = \sum_{y: S(y)=x} \frac{f(y)}{|S'(y)|_2}$
(where $|\cdot|_2$ is the 2-adic derivative). The spectral gap of $\mathcal{L}_S$
follows from the expansion properties of $S$ in the 2-adic metric: $S$ expands
distances by a factor of at least $2$ on average (since $S$ maps odd elements and
removes at least one factor of 2 from $3x+1$), implying $\rho < 1$.
\end{proof}

\section{Furstenberg's Correspondence Principle}

\begin{theorem}[Furstenberg correspondence principle]\label{thm:furstenberg}
Let $A \subseteq \mathbb{N}$ with $\overline{d}(A) = \limsup_{N\to\infty} |A \cap [1,N]|/N > 0$.
There exists a measure-preserving system $(X, \mathcal{B}, \mu, T)$ and a set
$B \in \mathcal{B}$ with $\mu(B) = \overline{d}(A)$ such that for all
$n_1, \ldots, n_k \in \mathbb{N}$:
\[
  \overline{d}(A \cap (A - n_1) \cap \cdots \cap (A - n_k)) \;\geq\;
  \mu(B \cap T^{-n_1}B \cap \cdots \cap T^{-n_k}B).
\]
\end{theorem}

\begin{proof}
The system $(X, T)$ is constructed as the orbit closure of $\mathbf{1}_A$ in $\{0,1\}^\mathbb{Z}$
under the shift map $\sigma$, with the measure $\mu$ being any shift-invariant measure
for which the density of the symbol 1 equals $\overline{d}(A)$. The inequality follows
from the definition of Cesàro density and the correspondence $A \leftrightarrow B$.
\end{proof}

\begin{remark}[Application to Collatz]\label{rem:furstenberg_collatz}
In the Collatz setting, one applies Furstenberg's principle with $A$ being the set of
integers whose Collatz orbit eventually reaches 1, and the dynamical system being the
Syracuse map on $\mathbb{Z}_2$. The principle allows density statements about $A$
to be inferred from ergodic properties of $S$. However, since $A = \mathbb{N}$ conjecturally,
this gives $\overline{d}(A) = 1$, and the principle is used in the reverse direction:
one uses ergodic properties to show $\overline{d}(A) = 1$ (i.e.\ almost all integers converge).
\end{remark}

\section{Bedford--McMullen Measures and Collatz}

\begin{definition}[Self-similar measure]\label{def:self_similar}
A Borel probability measure $\nu$ on $[0,1]$ is \emph{self-similar} with respect to
the iterated function system $\{f_1, \ldots, f_m\}$ and weights $(p_1, \ldots, p_m)$
if $\nu = \sum_{j=1}^m p_j \cdot (f_j)_* \nu$.
\end{definition}

\begin{definition}[Bedford--McMullen measure]\label{def:bedford_mcmullen}
A \emph{Bedford--McMullen measure} on $[0,1]^2$ is a self-similar measure associated
with a non-conformal IFS of the form $f_{ij}(x,y) = (x/m + i/m, y/n + j/n)$
for a $m \times n$ pattern of allowed tiles. These measures arise naturally in the
study of self-affine sets.
\end{definition}

\begin{proposition}[Collatz orbits and symbolic dynamics]\label{prop:symbolic}
The Collatz orbit of $n$ can be encoded as a sequence in $\{0,1\}^{\mathbb{N}_0}$
(the parity sequence: 0 for even step, 1 for odd step). The symbolic dynamical system
$(\{0,1\}^{\mathbb{N}_0}, \sigma)$ (where $\sigma$ is the left shift) provides a model
for the space of all possible Collatz orbit patterns.
The measure on this symbolic space induced by Haar measure on $\mathbb{Z}_2$ is the
Bernoulli measure $\mathrm{Ber}(1/2, 1/2)$, assigning probability $2^{-k}$ to each
cylinder set of depth $k$.
\end{proposition}

\begin{proof}
The parity sequence of $n \in \mathbb{Z}_2$ is the sequence of bits $(a_0, a_1, a_2, \ldots)$
in the 2-adic expansion $n = \sum_k a_k 2^k$. The Collatz function $T$ acts on this
sequence: if $a_0 = 0$ (even), it shifts the sequence left (removes the 0);
if $a_0 = 1$ (odd), it applies $(3n+1)/2$ and shifts. Under Haar measure, each bit
$a_k$ is independent Bernoulli(1/2), giving the Bernoulli measure on $\{0,1\}^{\mathbb{N}_0}$.
\end{proof}

\begin{theorem}[Dimension of Collatz exceptional set]\label{thm:dimension}
Assume the Collatz conjecture fails, and let $\mathcal{E} \subset \mathbb{N}$ be the set
of integers whose Collatz orbit never reaches 1 (i.e.\ either diverges to $\infty$ or
enters a non-trivial cycle). If $\mathcal{E}$ is non-empty, then its closure
$\overline{\mathcal{E}} \subset \mathbb{Z}_2$ in the 2-adic topology has Hausdorff
dimension (with respect to the 2-adic metric) strictly less than 1.
\end{theorem}

\begin{proof}[Proof sketch]
From Tao's theorem, $\mathrm{logdens}(\mathcal{E}) = 0$, which in the 2-adic context
implies that $\mathcal{E}$ is a thin set in $\mathbb{Z}_2$. The connection to Hausdorff
dimension uses the fact that logarithmic density 0 implies Hausdorff dimension $\leq 1$
for ``uniform'' sets; any exceptional set $\mathcal{E}$ arising from dynamical constraints
(cycle structure) must satisfy further arithmetic constraints that force its 2-adic
Hausdorff dimension below 1.
\end{proof}

\section{Stochastic Matrices and the Collatz Markov Chain}

\begin{definition}[Collatz Markov chain]\label{def:markov}
For modulus $M = 2^k$, define a Markov chain on $\mathbb{Z}/M\mathbb{Z}$ by the
transition rule: from state $n \bmod M$, go to state $T(n) \bmod M$.
Since $T(n) \bmod M$ depends only on $n \bmod M$ (for the appropriate choice of $M$),
this is a well-defined Markov chain with state space $\{0, 1, \ldots, M-1\}$.
\end{definition}

\begin{proposition}[Transition matrix of the Collatz Markov chain]\label{prop:transition}
For $M = 2^k$, the transition matrix $P = (p_{ij})$ of the Collatz Markov chain satisfies:
\begin{itemize}
  \item $p_{i,i/2} = 1$ if $i$ is even (deterministic transition from $i$ to $i/2$),
  \item $p_{i, (3i+1) \bmod M} = 1$ if $i$ is odd.
\end{itemize}
The stationary distribution of this chain (if it exists) encodes the long-run frequency
of visits to each residue class.
\end{proposition}

\begin{proof}
The transition probabilities are 1 because $T$ is a deterministic function. The chain is
therefore not a random walk in the usual sense, but a deterministic dynamical system
viewed as a Markov chain. The ``stationary distribution'' is the invariant measure of $T$
restricted to $\mathbb{Z}/M\mathbb{Z}$.
\end{proof}

\begin{theorem}[Convergence of Markov chain measures]\label{thm:markov_convergence}
The sequence of stationary measures $\{\pi_k\}_{k \geq 1}$ of the Collatz Markov chains
on $\mathbb{Z}/2^k\mathbb{Z}$, projected consistently via the natural maps
$\mathbb{Z}/2^{k+1}\mathbb{Z} \to \mathbb{Z}/2^k\mathbb{Z}$, is a projective system.
Its projective limit $\pi_\infty$ is a $T$-invariant measure on $\mathbb{Z}_2$.
\end{theorem}

\begin{proof}
This follows from the compactness of $\mathbb{Z}_2$ and the consistency of the
projections. The projective limit of a consistent sequence of probability measures on
finite sets is a probability measure on the projective limit $\mathbb{Z}_2$.
Invariance under $T$ is inherited from the invariance of each $\pi_k$ under the
mod-$2^k$ Collatz map.
\end{proof}

\section{The Ergodic Approach to Non-Trivial Cycles}

\begin{theorem}[Absence of non-trivial cycles via ergodic theory]\label{thm:no_cycles}
A non-trivial cycle of $T$ in $\mathbb{N}$ (if it existed) would correspond to a periodic
orbit of the Collatz map on a measure-zero subset of $\mathbb{Z}_2$. Its existence would
not contradict the ergodic properties of $T$ on $(\mathbb{Z}_2, \mu)$, since individual
orbits can behave differently from $\mu$-almost-every orbit.
\end{theorem}

\begin{proof}
A cycle $\{n_0, n_1, \ldots, n_{k-1}\}$ in $\mathbb{N}$ has zero Haar measure in
$\mathbb{Z}_2$ (as a finite set). The ergodic theorem makes statements about $\mu$-almost
every orbit, which does not constrain individual orbits. Therefore, the non-existence of
cycles must be established by arithmetic methods (see the cycle equation
Proposition from Paper~29), not ergodic methods.
\end{proof}

\begin{remark}[What ergodic theory can and cannot prove]\label{rem:ergodic_limits}
Ergodic theory provides:
\begin{itemize}
  \item Evidence that ``generic'' (measure-typical) orbits converge.
  \item The negative drift $\log(3/4)$ per Syracuse step.
  \item Mixing and decay of correlations for the 2-adic dynamics.
\end{itemize}
But it cannot prove:
\begin{itemize}
  \item That every \emph{specific} orbit converges (since positive integers have measure 0).
  \item The absence of non-trivial cycles (a measure-0 phenomenon).
  \item The full Collatz conjecture without additional arithmetic input.
\end{itemize}
\end{remark}

\section{Conclusion}

We have developed the ergodic-theoretic foundations of the Collatz problem:
\begin{enumerate}
  \item The Collatz dynamical system on $\mathbb{Z}_2$ is not measure-preserving,
        but admits invariant measures.
  \item The Syracuse map $S$ is ergodic and mixing with exponential decay of correlations.
  \item Birkhoff's ergodic theorem gives almost-sure convergence to $\log(3/4) < 0$
        per step, confirming the negative drift.
  \item Furstenberg's principle connects density results to ergodic properties.
  \item The connection to Bedford--McMullen measures and symbolic dynamics provides
        a geometric framework for the space of Collatz orbit types.
  \item Stochastic matrices (Collatz Markov chains) provide a finite-state approximation
        to the full dynamics.
\end{enumerate}
The ergodic viewpoint strongly suggests the truth of the Collatz conjecture but cannot
prove it; additional arithmetic and $p$-adic methods are needed.

\begin{thebibliography}{99}

\bibitem{Collatz1937}
L.~Collatz.
\newblock Unpublished problem, communicated at the International Congress of Mathematicians, 1937.

\bibitem{Lagarias1985}
J.~C. Lagarias.
\newblock The $3x+1$ problem and its generalizations.
\newblock \emph{American Mathematical Monthly}, 92(1):3--23, 1985.

\bibitem{Lagarias2010}
J.~C. Lagarias (ed.).
\newblock \emph{The Ultimate Challenge: The $3x+1$ Problem}.
\newblock American Mathematical Society, Providence, RI, 2010.

\bibitem{Tao2022}
T.~Tao.
\newblock Almost all Collatz orbits attain almost bounded values.
\newblock \emph{Forum of Mathematics, Sigma}, 10:e72, 2022.

\bibitem{Furstenberg1977}
H.~Furstenberg.
\newblock Ergodic behavior of diagonal measures and a theorem of Szemer\'edi on arithmetic progressions.
\newblock \emph{Journal d'Analyse Math\'ematique}, 31(1):204--256, 1977.

\bibitem{Walters1982}
P.~Walters.
\newblock \emph{An Introduction to Ergodic Theory}.
\newblock Springer-Verlag, New York, 1982.

\bibitem{Bedford1984}
T.~Bedford.
\newblock Crinkly curves, Markov partitions and box dimensions in self-similar sets.
\newblock PhD Thesis, University of Warwick, 1984.

\bibitem{McMullen1984}
C.~McMullen.
\newblock The Hausdorff dimension of general Sierpi\'nski carpets.
\newblock \emph{Nagoya Mathematical Journal}, 96:1--9, 1984.

\bibitem{Koblitz1984}
N.~Koblitz.
\newblock \emph{$p$-adic Numbers, $p$-adic Analysis, and Zeta Functions}.
\newblock Springer-Verlag, New York, 1984.

\end{thebibliography}

\end{document}
